\newlabel{C.Motivation}{{1}{7}{Motivating Examples}{chapter.1}{}}
\newlabel{fig:ch01:ch01_roadmap}{{1.1}{7}{The four parts of the course and the dependencies between the lectures (arrows). Generated by \texttt {code/ch01\_roadmap.py}}{figure.caption.1}{}}
\newlabel{fig:ch01:hardware}{{1.2}{9}{A superconducting quantum computer, from the outside in: the dilution cryostat, the wiring and filtering stages inside the can, the printed circuit board at its bottom (zoom), and the qubit chip at the centre of the board (zoom); to the right, the room-temperature electronics that generate the control pulses. Photographs: IBM; composition by \texttt {code/ch01\_hardware.py}}{figure.caption.2}{}}
\newlabel{fig:ch01:qubit_table}{{1.3}{10}{Physical realisations of qubits and their characteristic time scales. Source: \url {https://en.wikipedia.org/wiki/Qubit}}{figure.caption.3}{}}
\newlabel{fig:ch01:mckinsey}{{1.4}{10}{The industry view of the qubit platforms. Source: McKinsey, \emph {The year of quantum: from concept to reality in 2025}, \url {https://www.mckinsey.com/capabilities/tech-and-ai/our-insights/the-year-of-quantum-from-concept-to-reality-in-2025}}{figure.caption.4}{}}
\newlabel{fig:ch01:superconducting}{{1.5}{11}{Superconducting qubits: from the Josephson junction to the transmon. Source: \url {https://www.nature.com/articles/s41578-021-00370-4}}{figure.caption.5}{}}
\newlabel{tab:ch01:quantum_dots}{{1.1}{11}{Spin qubits in quantum dots: holes in strained germanium versus electrons in silicon. $^\dagger $The community accepts $20\,\mu $s. Source: \url {https://aip.scitation.org/doi/full/10.1063/1.5010933}}{table.caption.6}{}}
\newlabel{fig:ch01:quantum_dots}{{1.6}{11}{Double quantum dots. Source: \url {https://aip.scitation.org/doi/full/10.1063/1.5010933}}{figure.caption.7}{}}
\newlabel{eq:ch01:schrodinger}{{1.1}{11}{Quantum systems as controlled dynamical systems}{equation.1.1}{}}
\newlabel{fig:ch01:neutral_atoms}{{1.7}{12}{A two-dimensional array of neutral atoms in optical tweezers. Source: \url {https://www.nature.com/articles/s41586-022-04592-6}}{figure.caption.8}{}}
\newlabel{def:ch01:ic}{{1.1}{12}{Informationally complete measurement}{thm.1.1}{}}
\newlabel{eq:ch01:lininv}{{1.2}{12}{Example A: quantum state tomography}{equation.1.2}{}}
\newlabel{fig:ch01:pauli_bases}{{1.8}{13}{The three Pauli measurements of a qubit and their six eigenstates on the Bloch sphere. Their outcome statistics determine the Bloch vector, hence the state. Generated by \texttt {code/ch01\_pauli\_bases.py}}{figure.caption.9}{}}
\newlabel{fig:ch01:tomography_cloud}{{1.9}{13}{Linear-inversion estimates of the Bloch vector of a nearly pure qubit state from $N$ shots per Pauli axis ($150$ repetitions, $(r_x,r_z)$ plane). Red points lie outside the Bloch ball and are not density operators; the arrow shows the projection of one of them back onto the ball. Generated by \texttt {code/ch01\_tomography\_cloud.py}}{figure.caption.10}{}}
\newlabel{eq:ch01:pls}{{1.3}{13}{Example A: quantum state tomography}{equation.1.3}{}}
\newlabel{eq:ch01:mle}{{1.4}{13}{Example A: quantum state tomography}{equation.1.4}{}}
\newlabel{fig:ch01:tomography_shots}{{1.10}{14}{Mean trace distance between the estimate and the true qubit state as a function of the number of shots per axis, for a pure and for a mixed state, with and without projection onto the Bloch ball. The error follows the $1/\sqrt N$ law of Monte-Carlo estimation; the projection helps for the pure state and is neutral for the mixed one. Generated by \texttt {code/ch01\_tomography\_shots.py}}{figure.caption.11}{}}
\newlabel{fig:ch01:tomography_scaling}{{1.11}{14}{Number of real parameters of an $n$-qubit density operator and an order-of-magnitude count of the shots needed to estimate them to precision $10^{-2}$, against a budget of $10^9$ shots. Generated by \texttt {code/ch01\_tomography\_scaling.py}}{figure.caption.12}{}}
\newlabel{eq:ch01:rabi}{{1.5}{14}{Example B: a driven qubit}{equation.1.5}{}}
\newlabel{prob:ch01:transfer}{{1.2}{14}{Time-optimal state transfer, one qubit}{thm.1.2}{}}
\newlabel{fig:ch01:rabi}{{1.12}{15}{Left: excited-state population of the qubit \eqref {eq:ch01:rabi} under a constant drive $u=0.5$, on resonance (solid) and detuned by $\omega =2u$ (dashed). Right: the corresponding Bloch-sphere trajectories, from the north pole $\ket 0$: on resonance, a great circle to the south pole $\ket 1$; detuned, a smaller circle that never reaches $\ket 1$. Generated by \texttt {code/ch01\_rabi.py}}{figure.caption.13}{}}
\newlabel{fig:ch01:bangbang}{{1.13}{15}{A resonant pulse and a bang-bang pulse of bounded amplitude $u_0$ for $H=\tfrac {\omega }{2}\sigma _z+u(t)\sigma _x$ (left), and the population of $\ket 1$ they produce (right). With a bound on $u$, the bang-bang pulse is the shape that the Pontryagin maximum principle of \Chap ~\ref {C.GeometryControl} predicts. Generated by \texttt {code/ch01\_bangbang.py}}{figure.caption.14}{}}
\newlabel{eq:ch01:fixedtime}{{1.6}{15}{Example C: synthesising a gate in a landscape full of traps}{equation.1.6}{}}
\newlabel{fig:ch01:landscape}{{1.14}{16}{Fidelity $F=\abs {\Tr (U^{\star \dagger }U(u_1,u_2))}/3$ of a two-segment piecewise-constant control for the three-level system $\Hdrift =\operatorname {diag}(0,0.516,1)$, $V_{12}=1/\sqrt 2$, $V_{23}=1$ of \cite {bondar_uncomputability_2020}, for a target reachable at $(u_1,u_2)\approx (-0.33,0.88)$. The landscape has many ridges and local maxima. Right: the ten lines of Python that compute $U(T)$ and $F$. Generated by \texttt {code/ch01\_landscape.py}}{figure.caption.15}{}}
\newlabel{eq:ch01:qcpop}{{1.7}{16}{Example C: synthesising a gate in a landscape full of traps}{equation.1.7}{}}
\newlabel{eq:ch01:hamid}{{1.8}{16}{Example C: synthesising a gate in a landscape full of traps}{equation.1.8}{}}
\newlabel{fig:ch01:ch01_trap_slice}{{1.15}{17}{The fidelity of the three-level transfer along the straight segment from a local maximum of the landscape (a trap) to the global optimum. A gradient method started near the trap stays there. Generated by \texttt {code/ch01\_trap\_slice.py}}{figure.caption.16}{}}
\newlabel{eq:ch01:rse}{{1.9}{17}{Example D: a qubit with random noise}{equation.1.9}{}}
\newlabel{fig:ch01:ch01_rse_paths}{{1.16}{17}{Five sample paths of $\Tr (\rho \sigma _z)$ under the random Schr\"odinger equation of a driven qubit, each unitary, together with the average state and its purity: the average is a mixed state, and the randomness plays the role of decoherence. Generated by \texttt {code/ch01\_rse\_paths.py}}{figure.caption.17}{}}
\newlabel{eq:ch01:nonoisecost}{{1.10}{18}{Example D: a qubit with random noise}{equation.1.10}{}}
\newlabel{eq:ch01:noisecost}{{1.11}{18}{Example D: a qubit with random noise}{equation.1.11}{}}
\newlabel{eq:ch01:twoqubit}{{1.12}{18}{Example E: identifying a two-qubit open system}{equation.1.12}{}}
\newlabel{eq:ch01:twoqubitH}{{1.13}{18}{Example E: identifying a two-qubit open system}{equation.1.13}{}}
\newlabel{fig:ch01:ch01_twoqubit_data}{{1.17}{18}{Three Pauli expectation values of the two-qubit open system in time (curves) and their estimates from $200$ shots per time point (dots). Identification means recovering the generator of the dynamics from such data. Generated by \texttt {code/ch01\_twoqubit\_data.py}}{figure.caption.18}{}}
\newlabel{eq:ch01:hard}{{1.14}{19}{Example F: gates from local interactions}{equation.1.14}{}}
\newlabel{eq:ch01:soft}{{1.15}{19}{Example F: gates from local interactions}{equation.1.15}{}}
\newlabel{fig:ch01:ch01_penalty}{{1.18}{19}{Minimum-energy synthesis of a target gate on a qubit when rotations about $y$ are penalised by a factor $q$: as $q$ grows, the penalised direction is used less and less, and the control approaches one that only uses the allowed direction and its commutators. Generated by \texttt {code/ch01\_penalty.py}}{figure.caption.19}{}}
\newlabel{eq:ch01:conic}{{1.16}{20}{Conic optimization: a refresher}{equation.1.16}{}}
\newlabel{fig:ch01:cones}{{1.19}{21}{The nonnegative orthant, the Lorentz cone, and a spectrahedron: the elliptope of $3\times 3$ correlation matrices $\begin {psmallmatrix}1&x&y\\x&1&z\\y&z&1\end {psmallmatrix}\succeq 0$, an affine slice of the positive semidefinite cone. Its boundary is the cubic surface $1+2xyz-x^2-y^2-z^2=0$, curved everywhere except at the four vertices (the rank-one matrices); this is what distinguishes semidefinite from linear programming. Generated from numpy meshes by \texttt {code/ch01\_cones.py}}{figure.caption.20}{}}
\newlabel{fig:ch01:ch01_spectrahedron}{{1.20}{22}{An affine slice of the $3\times 3$ positive semidefinite cone (a spectrahedron) and the box cut out by its $2\times 2$ minors. Feasible sets of semidefinite programs look like this. Generated by \texttt {code/ch01\_spectrahedron.py}}{figure.caption.21}{}}
\newlabel{eq:ch01:pd}{{1.17}{22}{Conic optimization: a refresher}{equation.1.17}{}}
\newlabel{thm:ch01:duality}{{1.3}{22}{Weak and strong duality}{thm.1.3}{}}
\newlabel{fig:ch01:ch01_duality}{{1.21}{23}{A conic program over a slice of the second-order cone: the feasible chord, the level sets of the objective, and the optimal level set touching the feasible set, which is what the dual certificate expresses. Generated by \texttt {code/ch01\_duality.py}}{figure.caption.22}{}}
\newlabel{eq:ch01:sos}{{1.18}{23}{Conic optimization: a refresher}{equation.1.18}{}}
\newlabel{fig:ch01:ch01_sos}{{1.22}{23}{A quartic $f$, its minimum $f^\star $, and the two squares with $f-f^\star =s_1^2+s_2^2$: the certificate that $f^\star $ is the global minimum. Generated by \texttt {code/ch01\_sos.py}}{figure.caption.23}{}}
\newlabel{eq:ch01:qm}{{1.19}{24}{Conic optimization: a refresher}{equation.1.19}{}}
\newlabel{thm:ch01:putinar}{{1.4}{24}{Putinar \cite {putinar1993positive}}{thm.1.4}{}}
\newlabel{eq:ch01:pop}{{1.20}{24}{Conic optimization: a refresher}{equation.1.20}{}}
\newlabel{eq:ch01:hierarchy}{{1.21}{24}{Conic optimization: a refresher}{equation.1.21}{}}
\newlabel{prop:ch01:monotone}{{1.6}{24}{Nesterov, Parrilo, Lasserre \cite {nesterov2000squared,lasserre2001global,parrilo2000structured}}{thm.1.6}{}}
\newlabel{thm:ch01:lasserre}{{1.7}{24}{Convergence}{thm.1.7}{}}
\newlabel{C.Closed}{{2}{25}{Models of Closed Quantum Systems}{chapter.2}{}}
\newlabel{def:ch02:povm}{{2.1}{25}{POVM}{thm.2.1}{}}
\newlabel{eq:schrodinger}{{2.1}{25}{The Schr\"odinger equation and the propagator}{equation.2.1}{}}
\newlabel{eq:propagator-eq}{{2.2}{25}{The Schr\"odinger equation and the propagator}{equation.2.2}{}}
\newlabel{prop:unitary}{{2.2}{25}{}{thm.2.2}{}}
\newlabel{eq:control-hamiltonian}{{2.3}{26}{Controlled Hamiltonians as bilinear systems}{equation.2.3}{}}
\newlabel{eq:bilinear-group}{{2.4}{26}{Controlled Hamiltonians as bilinear systems}{equation.2.4}{}}
\newlabel{eq:pwc}{{2.5}{26}{Controlled Hamiltonians as bilinear systems}{equation.2.5}{}}
\newlabel{def:ch02:lie}{{2.3}{26}{}{thm.2.3}{}}
\newlabel{def:dla}{{2.4}{26}{Dynamical Lie algebra}{thm.2.4}{}}
\newlabel{fig:ch02:pwc}{{2.1}{27}{A piecewise-constant control (step function) and the population $\abs {\braket {1|\psi (t)}}^2$ it produces on a qubit with drift $\tfrac 12\sigma _z$ and control Hamiltonian $\tfrac 12\sigma _x$. Generated by \texttt {code/ch02\_pwc.py}}{figure.caption.24}{}}
\newlabel{def:reach}{{2.5}{27}{}{thm.2.5}{}}
\newlabel{thm:js-preview}{{2.6}{27}{Jurdjevic--Sussmann; preview}{thm.2.6}{}}
\newlabel{def:density}{{2.7}{27}{}{thm.2.7}{}}
\newlabel{eq:lvn}{{2.6}{27}{Density operators}{equation.2.6}{}}
\newlabel{eq:lvn-vec}{{2.7}{27}{Density operators}{equation.2.7}{}}
\newlabel{eq:ch02:bloch-rho}{{2.8}{28}{The qubit: Bloch equations}{equation.2.8}{}}
\newlabel{eq:ch02:qubitH}{{2.9}{28}{The qubit: Bloch equations}{equation.2.9}{}}
\newlabel{eq:ch02:bloch}{{2.10}{28}{The qubit: Bloch equations}{equation.2.10}{}}
\newlabel{fig:ch02:bloch}{{2.2}{28}{Bloch-sphere trajectories from $\ket 0$ under $H=\tfrac \Delta 2\sigz +\tfrac 12\sigx $: resonant ($\Delta =0$, a $\pi $-pulse reaching $\ket 1$) and detuned ($\Delta =1$). Generated by \texttt {code/ch02\_bloch.py}}{figure.caption.25}{}}
\newlabel{eq:ch02:basis_conditions}{{2.11}{29}{The coherence vector in the $\liesu [N]$ basis}{equation.2.11}{}}
\newlabel{eq:ch02:decomposition_hamiltonian}{{2.12}{29}{The coherence vector in the $\liesu [N]$ basis}{equation.2.12}{}}
\newlabel{eq:ch02:coherence}{{2.13}{29}{The coherence vector in the $\liesu [N]$ basis}{equation.2.13}{}}
\newlabel{eq:ch02:antisymmetric_structure_const}{{2.14}{29}{The coherence vector in the $\liesu [N]$ basis}{equation.2.14}{}}
\newlabel{eq:ch02:symmetric_structure_const}{{2.15}{29}{The coherence vector in the $\liesu [N]$ basis}{equation.2.15}{}}
\newlabel{prop:ch02:lvn-coherence}{{2.8}{29}{Liouville--von Neumann in coherence-vector form}{thm.2.8}{}}
\newlabel{prop:ch02:su2q}{{2.9}{29}{}{thm.2.9}{}}
\newlabel{eq:ch02:fixedtime}{{2.16}{30}{Quantum control problems}{equation.2.16}{}}
\newlabel{exe:su4}{{2.10}{30}{}{thm.2.10}{}}
\newlabel{C.Open}{{3}{31}{Models of Open Quantum Systems}{chapter.3}{}}
\newlabel{eq:ch03:hamil_full}{{3.1}{31}{Open quantum systems}{equation.3.1}{}}
\newlabel{eq:ch03:assump_sep}{{3.2}{31}{Open quantum systems}{equation.3.2}{}}
\newlabel{eq:ch03:first_master}{{3.3}{31}{Open quantum systems}{equation.3.3}{}}
\newlabel{eq:ch03:cmn}{{3.5}{32}{The GKSL master equation}{equation.3.5}{}}
\newlabel{thm:ch03:gksl}{{3.1}{32}{Canonical GKSL form \cite {Hall2014,gorini1976completely}}{thm.3.1}{}}
\newlabel{eq:ch03:gksl}{{3.6}{32}{Canonical GKSL form \cite {Hall2014,gorini1976completely}}{equation.3.6}{}}
\newlabel{eq:ch03:lindblad}{{3.7}{32}{The GKSL master equation}{equation.3.7}{}}
\newlabel{eq:ch03:bloch-open}{{3.8}{32}{The GKSL master equation}{equation.3.8}{}}
\newlabel{fig:ch03:decoherence}{{3.1}{33}{Bloch components of a qubit with $\omega =2$, $\gamma _1=0.4$, $\gamma _\phi =0.3$, starting in $\ket +$: damped precession of $x,y$ on the timescale $T_2$ and relaxation of $z$ to the ground state on the timescale $T_1$. Generated by \texttt {code/ch03\_decoherence.py}}{figure.caption.26}{}}
\newlabel{fig:ch03:shrink}{{3.2}{33}{Three initial pure states under the same dissipative dynamics as in Figure~\ref {fig:ch03:decoherence}: every trajectory leaves the sphere, spirals inwards and converges to the ground state at the north pole. Generated by \texttt {code/ch03\_decoherence.py}}{figure.caption.27}{}}
\newlabel{eq:ch03:vec-lind}{{3.9}{33}{The GKSL master equation}{equation.3.9}{}}
\newlabel{eq:ch03:nzme}{{3.10}{34}{Perturbative master equations}{equation.3.10}{}}
\newlabel{eq:ch03:born}{{3.11}{34}{Perturbative master equations}{equation.3.11}{}}
\newlabel{eq:ch03:redfield}{{3.12}{34}{Perturbative master equations}{equation.3.12}{}}
\newlabel{eq:ch03:secular}{{3.14}{35}{Perturbative master equations}{equation.3.14}{}}
\newlabel{sec:ch03:dynrep}{{4}{36}{The GKSL equation as a bilinear system}{section.3.4}{}}
\newlabel{eq:ch03:reformulated_control_hamiltonian}{{3.16}{36}{The GKSL equation as a bilinear system}{equation.3.16}{}}
\newlabel{thm:ch03:blds}{{3.2}{36}{Bilinear form of the controlled GKSL equation \cite {Parvaiz2025}}{thm.3.2}{}}
\newlabel{eq:ch03:BLDS_lindbladian}{{3.17}{36}{Bilinear form of the controlled GKSL equation \cite {Parvaiz2025}}{equation.3.17}{}}
\newlabel{eq:ch03:LDS_matrices_and_vectors}{{3.18}{36}{Bilinear form of the controlled GKSL equation \cite {Parvaiz2025}}{equation.3.18}{}}
\newlabel{eq:ch03:standard_form}{{3.19}{37}{The GKSL equation as a bilinear system}{equation.3.19}{}}
\newlabel{eq:ch03:output}{{3.20}{37}{The GKSL equation as a bilinear system}{equation.3.20}{}}
\newlabel{eq:ch03:accessible_set}{{3.21}{37}{The GKSL equation as a bilinear system}{equation.3.21}{}}
\newlabel{def:ch03:rode}{{3.3}{37}{RODE}{thm.3.3}{}}
\newlabel{eq:ch03:RODEmain}{{3.22}{37}{RODE}{equation.3.22}{}}
\newlabel{eq:ch03:rse}{{3.23}{38}{The random Schr\"odinger equation}{equation.3.23}{}}
\newlabel{prop:ch03:mixed}{{3.4}{38}{Mixed-state behaviour}{thm.3.4}{}}
\newlabel{fig:ch03:rse}{{3.3}{38}{Left: three sample paths of $\langle \sigma _x\rangle $ under the random Schr\"odinger equation with Ornstein--Uhlenbeck dephasing noise, and the ensemble mean over $300$ paths. Right: purity of the averaged state. Generated by \texttt {code/ch03\_rse.py}}{figure.caption.28}{}}
\newlabel{C.StateSpace}{{4}{41}{Conditions for Observability and Controllability in State Space}{chapter.4}{}}
\newlabel{eq:ch04:LDS}{{4.1}{41}{Linear systems}{equation.4.1}{}}
\newlabel{def:ch04:controllability}{{4.1}{41}{Controllability}{thm.4.1}{}}
\newlabel{def:ch04:observability}{{4.2}{41}{Observability}{thm.4.2}{}}
\newlabel{eq:ch04:OMCM}{{4.2}{41}{Linear systems}{equation.4.2}{}}
\newlabel{thm:ch04:kalman}{{4.3}{42}{Kalman}{thm.4.3}{}}
\newlabel{eq:ch04:BDS}{{4.3}{42}{Bilinear systems}{equation.4.3}{}}
\newlabel{eq:ch04:BDS_reach_and_obs}{{4.4}{42}{Bilinear systems}{equation.4.4}{}}
\newlabel{eq:ch04:similarity}{{4.5}{42}{Bilinear systems}{equation.4.5}{}}
\newlabel{thm:ch04:similarity}{{4.4}{42}{Minimal realisations are unique up to similarity}{thm.4.4}{}}
\newlabel{thm:ch04:js}{{4.5}{43}{Jurdjevic--Sussmann}{thm.4.5}{}}
\newlabel{eq:ch04:affine_lie}{{4.6}{44}{Controllability of closed systems}{equation.4.6}{}}
\newlabel{eq:ch04:BLDS}{{4.7}{44}{Observability of open quantum systems}{equation.4.7}{}}
\newlabel{fig:ch04:observability}{{4.1}{45}{Dimension of the observable subspace $\linspan \{\mathbf C,\mathbf C\mathbf A,\dots ,\mathbf C\mathbf A^{k}\}$ of the (embedded, 16-dimensional) two-qubit model of Section~\ref {sec:ch04:example} for several readouts, computed by \texttt {code/ch04\_observability.py}. Only an informationally complete readout observes the full state}{figure.caption.29}{}}
\newlabel{eq:ch04:accessible_set}{{4.8}{45}{Observability of open quantum systems}{equation.4.8}{}}
\newlabel{sec:ch04:identifiability}{{5}{45}{Identifiability}{section.4.5}{}}
\newlabel{fig:ch04:singular}{{4.2}{46}{Singular values of the Kalman observability matrix of the embedded ($4$-dimensional) damped-qubit model for three choices of readout ($\omega =1$, $\gamma _1=0.2$, $\gamma _\phi =0.1$). Zero singular values are drawn at the lower axis limit}{figure.caption.30}{}}
\newlabel{def:ch04:standing}{{4.6}{46}{Standing assumptions}{thm.4.6}{}}
\newlabel{def:ch04:ident}{{4.7}{46}{Two notions of identifiability}{thm.4.7}{}}
\newlabel{eq:ch04:T1T2}{{4.9}{46}{Identifiability}{equation.4.9}{}}
\newlabel{eq:ch04:M}{{4.10}{46}{Identifiability}{equation.4.10}{}}
\newlabel{thm:ch04:general}{{4.8}{46}{General parameter reconstruction}{thm.4.8}{}}
\newlabel{eq:ch04:symmetric_a_d}{{4.11}{47}{Identifiability}{equation.4.11}{}}
\newlabel{thm:ch04:symmetric}{{4.9}{47}{Reconstruction for symmetric $\Koss $}{thm.4.9}{}}
\newlabel{prop:ch04:gauge}{{4.10}{48}{Gauge fixing}{thm.4.10}{}}
\newlabel{rmk:ch04:partial}{{4.11}{48}{Partial readout}{thm.4.11}{}}
\newlabel{thm:ch04:auto}{{4.12}{48}{Autonomous OQS: identifiability from multirate sampling}{thm.4.12}{}}
\newlabel{cor:ch04:autoclosed}{{4.13}{48}{Closed autonomous systems}{thm.4.13}{}}
\newlabel{eq:ch04:pulses}{{4.12}{48}{Identifiability}{equation.4.12}{}}
\newlabel{thm:ch04:driven}{{4.14}{48}{Driven OQS: identifiability under the input class $\mathcal V_\alpha $}{thm.4.14}{}}
\newlabel{sec:ch04:example}{{132}{49}{Identifiability}{framenumber.132}{}}
\newlabel{eq:ch04:twoqubit}{{4.13}{49}{Identifiability}{equation.4.13}{}}
\newlabel{eq:ch04:kossblock}{{4.14}{49}{Identifiability}{equation.4.14}{}}
\newlabel{C.Frequency}{{5}{51}{Conditions for Observability and Controllability in the Frequency Domain}{chapter.5}{}}
\newlabel{eq:ch05:lds}{{5.1}{51}{Transfer functions}{equation.5.1}{}}
\newlabel{eq:ch05:tf}{{5.2}{51}{Transfer functions}{equation.5.2}{}}
\newlabel{def:ch05:tf}{{5.1}{51}{Transfer function}{thm.5.1}{}}
\newlabel{eq:ch05:modes}{{5.3}{51}{Transfer functions}{equation.5.3}{}}
\newlabel{fig:ch05:timeseries}{{5.1}{52}{Expectation values $\langle \sigma _x\rangle (t)$, $\langle \sigma _z\rangle (t)$ of a driven damped qubit ($\omega =2\pi $, $u=1.5$, $\gamma _1=0.15$, $\gamma _\phi =0.1$), sampled every $\tau =0.05$, generated by \texttt {code/ch05\_spectrum.py}}{figure.caption.31}{}}
\newlabel{thm:ch05:hautus}{{5.2}{52}{Hautus}{thm.5.2}{}}
\newlabel{fig:ch05:spectrum}{{5.2}{53}{Magnitude of the discrete Fourier transform of $\langle \sigma _x\rangle (t)$ from Figure~\ref {fig:ch05:timeseries}. The dashed lines mark the imaginary parts of the eigenvalues of $\mathbf A$; the peak width reflects the real part}{figure.caption.32}{}}
\newlabel{thm:ch05:kalman}{{5.3}{53}{Kalman decomposition; minimality}{thm.5.3}{}}
\newlabel{eq:ch05:discrete}{{5.4}{53}{Markov parameters, Hankel matrices, and realisation}{equation.5.4}{}}
\newlabel{eq:ch05:hankel}{{5.5}{53}{Markov parameters, Hankel matrices, and realisation}{equation.5.5}{}}
\newlabel{thm:ch05:hankelrank}{{5.4}{53}{Hankel rank}{thm.5.4}{}}
\newlabel{eq:ch05:era}{{5.6}{54}{Markov parameters, Hankel matrices, and realisation}{equation.5.6}{}}
\newlabel{fig:ch05:poles}{{5.3}{55}{Eigenvalues (poles) of the damped-qubit generator $\mathbf A$ in the complex plane and the estimates $\log (\lambda _i[\mathbf A_d])/\tau $ obtained by Ho--Kalman from $400$ noiseless samples of $(\langle \sigma _x\rangle ,\langle \sigma _z\rangle )$}{figure.caption.33}{}}
\newlabel{eq:ch05:eigset}{{5.7}{55}{Sampling and aliasing}{equation.5.7}{}}
\newlabel{eq:ch05:intersection}{{5.8}{55}{Sampling and aliasing}{equation.5.8}{}}
\newlabel{eq:ch05:reconstructA}{{5.9}{55}{Sampling and aliasing}{equation.5.9}{}}
\newlabel{eq:ch05:pe}{{5.10}{56}{Identifiability in the frequency domain}{equation.5.10}{}}
\newlabel{C.SampleComplexity}{{6}{57}{Sample Complexity of Identification}{chapter.6}{}}
\newlabel{thm:ch06:hoeffding}{{6.1}{57}{Hoeffding}{thm.6.1}{}}
\newlabel{fig:ch06:shots}{{6.1}{58}{Root-mean-square error of the least-squares estimate of a decay rate $\gamma =0.5$ from binomial estimates of $\langle \sigma _z\rangle $ at $12$ times, as a function of the number of shots $N$ per time, over $200$ Monte-Carlo repetitions; the dashed line is $\propto N^{-1/2}$. Generated by \texttt {code/ch06\_shots.py}}{figure.caption.34}{}}
\newlabel{thm:ch06:tomography}{{6.2}{58}{Sample complexity of tomography; \citealp {haah2016sample}, O'Donnell--Wright}{thm.6.2}{}}
\newlabel{thm:ch06:shadows}{{6.3}{59}{\citealp {Huang2020}}{thm.6.3}{}}
\newlabel{def:ch06:fisher}{{6.4}{59}{Fisher information}{thm.6.4}{}}
\newlabel{thm:ch06:cramerrao}{{6.5}{59}{Cram\'er--Rao; quantum Cram\'er--Rao}{thm.6.5}{}}
\newlabel{eq:ch06:ls}{{6.1}{60}{Finite-sample identification of linear systems}{equation.6.1}{}}
\newlabel{prop:ch06:hokalman}{{6.6}{60}{Shots for Ho--Kalman on quantum data}{thm.6.6}{}}
\newlabel{eq:ch06:error_bound}{{6.2}{60}{From the system matrix to the master equation}{equation.6.2}{}}
\newlabel{C.GeometryStates}{{7}{63}{Geometry of Quantum States}{chapter.7}{}}
\newlabel{eq:ch07:basis}{{7.1}{63}{Coherence vectors}{equation.7.1}{}}
\newlabel{eq:ch07:structure}{{7.2}{63}{Coherence vectors}{equation.7.2}{}}
\newlabel{eq:ch07:f}{{7.2a}{63}{Coherence vectors}{equation.7.2a}{}}
\newlabel{eq:ch07:g}{{7.2b}{63}{Coherence vectors}{equation.7.2b}{}}
\newlabel{def:ch07:coherence}{{7.1}{63}{Coherence vector}{thm.7.1}{}}
\newlabel{eq:ch07:coherence}{{7.3}{63}{Coherence vector}{equation.7.3}{}}
\newlabel{eq:ch07:su4}{{7.4}{64}{Coherence vectors}{equation.7.4}{}}
\newlabel{prop:ch07:sparsity}{{7.2}{64}{Sparsity of the structure constants}{thm.7.2}{}}
\newlabel{def:ch07:dens}{{7.3}{64}{}{thm.7.3}{}}
\newlabel{def:ch07:major}{{7.4}{64}{Majorisation}{thm.7.4}{}}
\newlabel{thm:ch07:uhlmann-major}{{7.5}{64}{Uhlmann}{thm.7.5}{}}
\newlabel{fig:ch07:blochbody}{{7.1}{65}{Planar sections of the state space in coherence coordinates, computed by bisection on the radius in each direction (\texttt {code/ch07\_blochbody.py}). Left: the qubit, a disc of radius one. Middle and right: two sections of the qutrit Bloch body in the Gell-Mann basis; the dotted and dashed circles are the inscribed and circumscribed spheres, $r_{\mathrm {in}}=1/\sqrt 3$ and $r_{\mathrm {out}}=2/\sqrt 3$. The $(x_3,x_8)$ section contains the diagonal states and is the probability simplex}{figure.caption.35}{}}
\newlabel{def:ch07:trace}{{7.6}{65}{}{thm.7.6}{}}
\newlabel{thm:ch07:helstrom}{{7.7}{65}{Helstrom}{thm.7.7}{}}
\newlabel{def:ch07:fidelity}{{7.8}{65}{Uhlmann fidelity}{thm.7.8}{}}
\newlabel{thm:ch07:fvdg}{{7.9}{65}{Uhlmann; Fuchs--van de Graaf}{thm.7.9}{}}
\newlabel{eq:ch07:fvdg}{{7.5}{65}{Uhlmann; Fuchs--van de Graaf}{equation.7.5}{}}
\newlabel{fig:ch07:metrics}{{7.2}{66}{Distances between $\rho _0=\ket 0\bra 0$ and $\rho (s)$ with Bloch vector $(0,0,1-2s)$, from \texttt {code/ch07\_metrics.py}. The trace and Hilbert--Schmidt distances are linear in the Bloch vector; the Bures angle has infinite slope at the pure state and saturates at $\pi /2$ once the states become orthogonal, at $s=1$ where $\Fid =0$}{figure.caption.36}{}}
\newlabel{eq:ch07:bures}{{7.6}{66}{Distances between states}{equation.7.6}{}}
\newlabel{eq:ch07:qcrb}{{7.7}{67}{Distances between states}{equation.7.7}{}}
\newlabel{C.GeometryChannels}{{8}{69}{Geometry of Quantum Channels}{chapter.8}{}}
\newlabel{eq:ch08:dynmap}{{8.1}{69}{Dynamical maps}{equation.8.1}{}}
\newlabel{eq:ch08:trace}{{8.2}{69}{Dynamical maps}{equation.8.2}{}}
\newlabel{eq:ch08:cp}{{8.3}{69}{Dynamical maps}{equation.8.3}{}}
\newlabel{eq:ch08:kraus}{{8.4}{69}{Dynamical maps}{equation.8.4}{}}
\newlabel{thm:ch08:choi}{{8.1}{70}{Choi~\citeyearpar {CHOI1975285}}{thm.8.1}{}}
\newlabel{eq:ch08:slice}{{8.5}{70}{The convex body of channels}{equation.8.5}{}}
\newlabel{thm:ch08:extreme}{{8.2}{70}{Choi; extreme points}{thm.8.2}{}}
\newlabel{fig:ch08:channels}{{8.1}{71}{The $(x,z)$-section of the Bloch ball (black) and its image under amplitude damping with $\gamma \in \{0.3,0.6,0.9\}$ (left; the image shifts towards $\ket 0$ at the top) and under dephasing with $p\in \{0.15,0.3,0.45\}$ (right; the $z$-axis is fixed, coherences shrink by $1-2p$). Computed from the Kraus operators in \texttt {code/ch08\_channels.py}, which also verifies CP and TP through the Choi matrix}{figure.caption.37}{}}
\newlabel{eq:ch08:classical-markov}{{8.6}{71}{Quantum Markovianity and divisibility}{equation.8.6}{}}
\newlabel{eq:ch08:composition}{{8.7}{71}{Quantum Markovianity and divisibility}{equation.8.7}{}}
\newlabel{def:ch08:divisible}{{8.3}{72}{P- and CP-divisibility}{thm.8.3}{}}
\newlabel{prop:ch08:cpdiv}{{8.4}{72}{Consequence of CP-divisibility~\cite {wolf06,Rivas2011OpenQS,Kossakowski10}}{thm.8.4}{}}
\newlabel{eq:ch08:timelocal}{{8.8}{72}{Consequence of CP-divisibility~\cite {wolf06,Rivas2011OpenQS,Kossakowski10}}{equation.8.8}{}}
\newlabel{fig:ch08:divisibility}{{8.2}{72}{Trace distance between two evolving qubit states under pure dephasing with coherence factor $c(t)$, from \texttt {code/ch08\_divisibility.py} ($\gamma =0.4$, $\omega =2$). Monotone decay for the semigroup; revivals, i.e.\ information flowing back from the environment, for the non-Markovian toy model}{figure.caption.38}{}}
\newlabel{eq:ch08:semigroup}{{8.9}{72}{Quantum Markovianity and divisibility}{equation.8.9}{}}
\newlabel{fig:ch08:hierarchy}{{8.3}{73}{Nested classes of trace-preserving maps, from Hermiticity-preserving linear maps down to unitary conjugations; the Markovian descriptions of \Chap ~\ref {C.Open} live in the two innermost classes. Drawn by \texttt {code/ch08\_hierarchy.py}}{figure.caption.39}{}}
\newlabel{def:ch08:distances}{{8.5}{73}{}{thm.8.5}{}}
\newlabel{C.GeometryControl}{{9}{75}{Geometry of Quantum Optimal Control}{chapter.9}{}}
\newlabel{eq:ch09:dyn}{{9.1}{75}{Optimal control problems}{equation.9.1}{}}
\newlabel{eq:ch09:cost}{{9.2}{75}{Optimal control problems}{equation.9.2}{}}
\newlabel{def:ch09:action}{{9.1}{75}{Action of an optimal control problem}{thm.9.1}{}}
\newlabel{eq:ch09:action}{{9.3}{75}{Action of an optimal control problem}{equation.9.3}{}}
\newlabel{thm:ch09:extremals}{{9.2}{76}{Extremals of the action}{thm.9.2}{}}
\newlabel{eq:ch09:el1}{{9.4}{76}{Extremals of the action}{equation.9.4}{}}
\newlabel{eq:ch09:el2}{{9.5}{76}{Extremals of the action}{equation.9.5}{}}
\newlabel{eq:ch09:el3}{{9.6}{76}{Extremals of the action}{equation.9.6}{}}
\newlabel{def:ch09:HP}{{9.3}{76}{Pontryagin Hamiltonian}{thm.9.3}{}}
\newlabel{eq:ch09:HP}{{9.7}{76}{Pontryagin Hamiltonian}{equation.9.7}{}}
\newlabel{thm:ch09:pmp}{{9.4}{76}{Pontryagin maximum principle}{thm.9.4}{}}
\newlabel{eq:ch09:hamilton}{{9.8}{76}{Pontryagin maximum principle}{equation.9.8}{}}
\newlabel{eq:ch09:max}{{9.9}{76}{Pontryagin maximum principle}{equation.9.9}{}}
\newlabel{exm:ch09:particle}{{9.5}{77}{Time-optimal control of a particle}{thm.9.5}{}}
\newlabel{eq:ch09:qlag}{{9.10}{77}{From the Schr\"odinger equation to the Pontryagin Hamiltonian}{equation.9.10}{}}
\newlabel{eq:ch09:qHP}{{9.11}{77}{From the Schr\"odinger equation to the Pontryagin Hamiltonian}{equation.9.11}{}}
\newlabel{eq:ch09:qadj}{{9.12}{77}{From the Schr\"odinger equation to the Pontryagin Hamiltonian}{equation.9.12}{}}
\newlabel{eq:ch09:qmax}{{9.13}{77}{From the Schr\"odinger equation to the Pontryagin Hamiltonian}{equation.9.13}{}}
\newlabel{eq:ch09:bloch}{{9.14}{77}{From the Schr\"odinger equation to the Pontryagin Hamiltonian}{equation.9.14}{}}
\newlabel{eq:ch09:pipulse}{{9.15}{77}{Time-optimal control of a qubit}{equation.9.15}{}}
\newlabel{eq:ch09:Hx}{{9.16}{78}{Time-optimal control of a qubit}{equation.9.16}{}}
\newlabel{eq:ch09:twobangs}{{9.17}{78}{Time-optimal control of a qubit}{equation.9.17}{}}
\newlabel{fig:ch09:timeoptimal}{{9.1}{78}{Left: Bloch-sphere trajectories for $\Delta =0.6$, $u_0=1$. The bang-bang solution \eqref {eq:ch09:twobangs} (blue) switches on the equator; the drift-free $\pi $-pulse (orange) is the pole-to-pole geodesic; a weak resonant drive $u=\tfrac 12u_0\cos \Delta t$ (green) spirals slowly towards the target. Right: the bang-bang control and the switching function $H_x(t)$ computed by integrating $\ket \psi $ and $\ket \chi $; $u$ flips sign where $H_x$ vanishes. Generated by \texttt {code/ch09\_timeoptimal.py}}{figure.caption.40}{}}
\newlabel{def:ch09:landscape}{{9.6}{79}{Control landscape}{thm.9.6}{}}
\newlabel{fig:ch09:landscape}{{9.2}{79}{Transfer fidelity $F(u_1,u_2)=\abs {\braket {\downarrow |U_2U_1|\uparrow }}^2$ for two piecewise-constant pulses of duration $\Delta t=1.6$ on a qubit with detuning $\Delta =0.6$. White dots mark local maxima, red crosses local minima on the grid. The landscape is periodic and, because the two-parameter family is far from covering the whole control space, it has several isolated maxima that are not global. Generated by \texttt {code/ch09\_landscape.py}}{figure.caption.41}{}}
\newlabel{eq:ch09:rse}{{9.18}{80}{Noise metrics and robust control}{equation.9.18}{}}
\newlabel{thm:ch09:basicerror}{{9.7}{80}{Basic error estimate}{thm.9.7}{}}
\newlabel{thm:ch09:geometricerror}{{9.8}{80}{Geometric error estimate}{thm.9.8}{}}
\newlabel{eq:ch09:bestbound}{{9.19}{80}{Geometric error estimate}{equation.9.19}{}}
\newlabel{def:ch09:wcnc}{{9.9}{81}{Worst-case noise condition, WCNC}{thm.9.9}{}}
\newlabel{thm:ch09:tightness}{{9.10}{81}{Tightness}{thm.9.10}{}}
\newlabel{eq:ch09:noisemetric}{{9.20}{81}{Noise metrics and robust control}{equation.9.20}{}}
\newlabel{eq:ch09:robust}{{9.21}{81}{Noise metrics and robust control}{equation.9.21}{}}
\newlabel{def:ch09:wcnc3}{{9.11}{81}{Control-dependent WCNC}{thm.9.11}{}}
\newlabel{thm:ch09:geodesic}{{9.12}{81}{Robust control as geodesics}{thm.9.12}{}}
\newlabel{C.Complexity}{{10}{83}{Quantum Circuit Complexity and the Length of the Trajectory in Quantum Optimal Control}{chapter.10}{}}
\newlabel{def:ch10:complexity}{{10.1}{83}{Circuit complexity}{thm.10.1}{}}
\newlabel{eq:ch10:nielsen}{{10.1}{83}{Circuit complexity and Nielsen's geometry}{equation.10.1}{}}
\newlabel{eq:ch10:lowerbound}{{10.2}{83}{Circuit complexity and Nielsen's geometry}{equation.10.2}{}}
\newlabel{fig:ch10:complexity}{{10.1}{84}{Left: unit balls of the penalised metrics $\gq $ in a plane spanned by a cheap (horizontal, in $D$) and an expensive (vertical, in $D^\perp $) direction; as $q\to \infty $ the ball flattens to the horizontal segment of the sub-Riemannian geometry. Right: a round geodesic from $\Id $ to $V$ spends most of its length in the expensive direction, so its $\gq $-length grows like $\sqrt q$; a horizontal path, which reaches the same target by combining cheap directions whose bracket generates $D^\perp $, has a $\gq $-length independent of $q$. Generated by \texttt {code/ch10\_complexity.py}}{figure.caption.42}{}}
\newlabel{eq:ch10:gq}{{10.3}{84}{Hard and soft constraints}{equation.10.3}{}}
\newlabel{eq:ch10:Eq}{{10.4}{84}{Hard and soft constraints}{equation.10.4}{}}
\newlabel{def:ch10:gamma}{{10.2}{85}{$\Gamma $-convergence}{thm.10.2}{}}
\newlabel{thm:ch10:fundamental}{{10.3}{85}{Fundamental theorem of $\Gamma $-convergence; \cite {DeGiorgi,dalMaso}}{thm.10.3}{}}
\newlabel{def:ch10:sobolev}{{10.4}{85}{Sobolev curves \cite {Lewis2020}}{thm.10.4}{}}
\newlabel{thm:ch10:gammaC0}{{10.5}{85}{$\Gamma $-convergence on $\Cxy $}{thm.10.5}{}}
\newlabel{thm:ch10:penalised}{{10.6}{85}{Convergence of penalised geodesics \cite {subRiemannian}}{thm.10.6}{}}
\newlabel{prop:ch10:full}{{10.7}{86}{Full sequence}{thm.10.7}{}}
\newlabel{prop:ch10:apriori}{{10.8}{86}{A priori bound on the constraint violation}{thm.10.8}{}}
\newlabel{eq:ch10:apriori}{{10.5}{86}{A priori bound on the constraint violation}{equation.10.5}{}}
\newlabel{eq:ch10:drift}{{10.6}{86}{Drift: the geometric interaction picture}{equation.10.6}{}}
\newlabel{def:ch10:flow}{{10.9}{86}{Assumption on the flow}{thm.10.9}{}}
\newlabel{prop:ch10:interaction}{{10.10}{86}{Geometric interaction picture}{thm.10.10}{}}
\newlabel{def:ch10:symmetry}{{10.11}{86}{Symmetry of the drift}{thm.10.11}{}}
\newlabel{lem:ch10:conformal}{{10.12}{86}{}{thm.10.12}{}}
\newlabel{prop:ch10:correspondence}{{10.13}{86}{Geodesic correspondence}{thm.10.13}{}}
\newlabel{thm:ch10:control}{{10.14}{87}{Convergence of optimal controls \cite {subRiemannian}}{thm.10.14}{}}
\newlabel{eq:ch10:qsystem}{{10.7}{87}{Gate synthesis with 2-local controls}{equation.10.7}{}}
\newlabel{fig:ch10:penalised}{{10.2}{88}{Penalised geodesics on $\SU [2]$ computed by minimising the discretised energy with a stiff endpoint penalty (24 piecewise-constant steps). Top left: Bloch trajectories of $U(t)\ket 0$ for $q=1$ (direct rotation about $y$), $q=16$ and $q=256$ (horizontal route). Top right: the penalised control component $a_y(t)$. Bottom: $\min \Ener {q}$ increasing to $\min \Ener {\infty }$, and $q\int g(\Pperp \dot \gamma _q,\Pperp \dot \gamma _q)$ tending to zero after the crossover. Toy computation, \texttt {code/ch10\_penalised.py}}{figure.caption.43}{}}
\newlabel{C.AlgTomography}{{11}{91}{Algorithmic Quantum State Tomography and Quantum Channel Tomography}{chapter.11}{}}
\newlabel{def:ch11:lininv}{{11.1}{91}{Linear inversion}{thm.11.1}{}}
\newlabel{def:ch11:mle}{{11.2}{92}{Maximum-likelihood estimate}{thm.11.2}{}}
\newlabel{prop:ch11:pls}{{11.3}{92}{Projected least squares}{thm.11.3}{}}
\newlabel{fig:ch11:tomography}{{11.1}{93}{Mean infidelity $1-F(\psi ,\hat \rho )$ of two tomographic estimators of a fixed pure qubit state, averaged over 400 Monte-Carlo runs, as a function of the number of shots $N$ per Pauli setting ($X$, $Y$, $Z$). Both scale as $1/N$ for large $N$; the projection onto the Bloch ball removes the unphysical excursions of linear inversion. Generated by \texttt {code/ch11\_tomography.py}}{figure.caption.44}{}}
\newlabel{eq:ch11:era}{{11.1}{94}{Identification of open systems from time series}{equation.11.1}{}}
\newlabel{eq:ch11:pop}{{11.2}{95}{Identification of open systems from time series}{equation.11.2}{}}
\newlabel{eq:ch11:T1T2}{{11.3}{95}{Recovery of master-equation parameters}{equation.11.3}{}}
\newlabel{eq:ch11:M}{{11.4}{95}{Recovery of master-equation parameters}{equation.11.4}{}}
\newlabel{thm:ch11:general}{{11.4}{95}{General reconstruction; \citealp {Parvaiz2025}}{thm.11.4}{}}
\newlabel{alg:ch11:general}{{11.1}{96}{Recovery of master-equation parameters}{algocf.11.1}{}}
\newlabel{alg:ch11:symmetric}{{11.2}{96}{Recovery of master-equation parameters}{algocf.11.2}{}}
\newlabel{eq:ch11:errorbound}{{11.5}{96}{Recovery of master-equation parameters}{equation.11.5}{}}
\newlabel{eq:ch11:gammablock}{{11.6}{96}{Recovery of master-equation parameters}{equation.11.6}{}}
\newlabel{fig:ch11:lindblad}{{11.2}{97}{Generator $\bA $ of the coherence-vector dynamics $\dot x=\bA x$ of a qubit ($\omega =2$, $\gamma _1=0.3$, $\gamma _\phi =0.2$), exact (left) and recovered by least squares from $4\times 80$ noisy time steps with $S=10^4$ shots per expectation value (right). The first row is zero (trace preservation), the first column carries the affine part $\vec \beta $ due to relaxation, the $XY$ block is the rotation at frequency $\omega $ damped at rate $\gamma _1/2+\gamma _\phi $. A toy re-computation in the spirit of the two-qubit experiments of \citet {Parvaiz2025}; generated by \texttt {code/ch11\_lindblad\_id.py}}{figure.caption.45}{}}
\newlabel{eq:ch11:Vunknown}{{11.7}{98}{Hamiltonian identification as polynomial optimization}{equation.11.7}{}}
\newlabel{C.AlgControl}{{12}{99}{Algorithmic Quantum Optimal Control}{chapter.12}{}}
\newlabel{eq:ch12:finite}{{12.1}{99}{What has to be computed}{equation.12.1}{}}
\newlabel{alg:ch12:shooting}{{12.1}{100}{Shooting methods}{algocf.12.1}{}}
\newlabel{eq:ch12:gd}{{12.2}{100}{Gradient methods tailored to quantum dynamics}{equation.12.2}{}}
\newlabel{eq:ch12:chain}{{12.3}{101}{Gradient methods tailored to quantum dynamics}{equation.12.3}{}}
\newlabel{eq:ch12:grape}{{12.4}{101}{Gradient methods tailored to quantum dynamics}{equation.12.4}{}}
\newlabel{fig:ch12:grape}{{12.1}{101}{GRAPE for the qubit transfer $\ket 0\to \ket 1$ with drift $\Delta =0.5$ and $N=40$ slices. Left: infidelity per iteration (log scale). Right: initial guess and optimised pulse. Generated by \texttt {code/ch12\_grape.py}}{figure.caption.46}{}}
\newlabel{prop:ch12:grape-pmp}{{12.1}{102}{}{thm.12.1}{}}
\newlabel{fig:ch12:traps}{{12.2}{102}{Objective $\langle O\rangle $ of the Pechen--Tannor problem as a function of the constant control $\lambda $ (drift $\operatorname {diag}(0,1,2)$). Crosses: local maxima; dot: global maximum. Generated by \texttt {code/ch12\_landscape\_traps.py}}{figure.caption.47}{}}
\newlabel{eq:ch12:qcc}{{12.5}{102}{Global optimisation: quantum control via polynomial optimisation}{equation.12.5}{}}
\newlabel{eq:ch12:robust-bound}{{12.6}{103}{Noise-informed control}{equation.12.6}{}}
\newlabel{eq:ch12:c2}{{12.7}{103}{Noise-informed control}{equation.12.7}{}}
\newlabel{C.Revisited}{{13}{105}{The Motivating Examples Revisited}{chapter.13}{}}
\newlabel{eq:ch13:example}{{13.1}{105}{Two coupled qubits: identifiability and identification}{equation.13.1}{}}
\newlabel{eq:ch13:gamma}{{13.2}{105}{Two coupled qubits: identifiability and identification}{equation.13.2}{}}
\newlabel{fig:ch13:twoqubit}{{13.1}{106}{Two-qubit model \eqref {eq:ch13:example}, qualitative re-creation. Left: relative errors of the Hamiltonian parameters and of the rates against the number of shots $S$ per time point, averaged over 5 runs; dashed: $S^{-1/2}$. Right: exact and recovered parameters at $S=10^4$. Generated by \texttt {code/ch13\_twoqubit\_id.py}}{figure.caption.48}{}}
\newlabel{fig:ch13:qcpop}{{13.2}{107}{Histogram of the final infidelities of 200 local gradient searches from random starts for a three-segment control of the three-level system; dashed: the global optimum found by enumeration. Generated by \texttt {code/ch13\_qcpop\_vs\_grape.py}}{figure.caption.49}{}}
\newlabel{eq:ch13:2local}{{13.3}{108}{2-local control of $n$ qubits: penalised geodesics}{equation.13.3}{}}
\newlabel{tocindent-1}{0pt}
\newlabel{tocindent0}{61.69456pt}
\newlabel{tocindent1}{22.77783pt}
\newlabel{tocindent2}{0pt}
\newlabel{tocindent3}{0pt}
